\documentclass{article}
\usepackage[margin=1in]{geometry}
\usepackage{energese}
\usepackage{hyperref}

\title{The \texttt{energese} Package User Guide}
\author{Jules}
\date{January 2026}

\begin{document}

\maketitle

\section{Introduction}
The \texttt{energese} package automates the rendering of H.T. Odum's Energy Systems Language (ESL) diagrams using a data-driven approach via LuaLaTeX and TikZ.

\section{Quick Start}
To use the package, you need to provide a JSON file describing your system.

\subsection{Example System}
Here is a simple ESL diagram rendered from a JSON file.

\begin{center}
\renderEnergese{examples/simple_system.json}
\end{center}

\subsection{JSON Structure}
The JSON file should contain \texttt{nodes} and \texttt{edges}. Each node must have a \texttt{Tr} (Transformity) value, which determines its horizontal position.

\begin{verbatim}
{
  "nodes": [
    {"id": "sun", "type": "source", "Tr": 1, "label": "Sun"},
    {"id": "grass", "type": "producer", "Tr": 100, "label": "Grass"}
  ],
  "edges": [
    {"from": "sun", "to": "grass", "type": "energy"}
  ]
}
\end{verbatim}

\section{Node Types}
The package supports several node types:
\begin{itemize}
    \item \texttt{source}: Circle
    \item \texttt{producer}: Lung shape
    \item \texttt{consumer}: Hexagon
    \item \texttt{storage}: Tank
    \item \texttt{interaction}: Pointed block
    \item \texttt{transaction}: Diamond
\end{itemize}

\section{Edge Types}
\begin{itemize}
    \item \texttt{energy}: Solid line with arrow
    \item \texttt{money\_feedback}: Dashed blue overhead arc
    \item \texttt{information}: Thin solid line
    \item \texttt{material}: Solid line
\end{itemize}

\end{document}
